\section{Introduction}
Consider the following scenario in university admissions. A university (the primary university) seeks to admit candidates with the right aptitude for its program. However, a candidate’s aptitude is not directly observable and must instead be inferred from their profile, which includes exam scores, experience, and other attributes. Typically, admission decisions are based on historical data, profiles of past candidates and their subsequent performance in the program.

Now, suppose a candidate, in addition to their own profile, also discloses the outcome of their application at another university (the secondary university). This additional information, referred to as \emph{side information}, raises a central question: How should the primary university leverage the decision of the secondary university to improve the accuracy of its own admission process?

The value of such side information depends crucially on the relation between the two universities. If the secondary university shares similar academic standards, subject focus, and reputation, the side information can be beneficial in assessing the candidate's aptitude. Conversely, if the two institutions differ significantly in their requirements or selectivity, the side information may be misleading and could result in suboptimal decisions. In general, the degree of alignment between the objectives and standards of the two universities determines how much weight the primary university should assign to the side information.

We consider the setting where, historically, no candidate has admission outcomes from both universities at the same time. Thus the past data includes rows with student features and admission decisions from either the primary university or the secondary university. Indeed if decisions of both universities were available, the secondary university decision could be absorbed readily as an additional feature in the learning problem. The data with the primary university decisions is also assumed to be more voluminous compared to that with secondary university decisions. Both these assumptions are motivated from real life experiences, students often apply to universities near their homes whereby past student data is often partitioned as above when the primary and secondary universities are geographically far apart. A university is also likely to have access to very limited data of other universities.

In this chapter, we formally study this problem and quantify how the error in admission decisions relates to the amount of trust placed in side information. We further characterize how much trust to place in side information, assuming the true relationship between candidate features, the primary university’s decision, and the secondary university’s decision is known a priori.

We formalize this problem within a classification framework. The primary and secondary university classifiers predict the admission outcome for a candidate, referred to as the label and side information, respectively.
The primary university trains its classifier using both its own training data and the secondary university's decision rule. This allows the primary university to incorporate the secondary university's knowledge into its own inference process.
At test time, when both the candidate’s profile and the secondary university’s admission decision are available, the primary university aims to construct a composite classifier that combines its own prediction with the side information, ensuring performance at least as good as the primary classifier alone.

Recall that the two training datasets are completely disjoint, referred to as a \textit{partitioned dataset}. The challenge with a partitioned dataset is that very different underlying relationships between candidate features, the primary university’s decision, and the secondary university’s decision can produce the same observed data. More precisely, the \textit{joint distribution} of features, label and side information is \textit{unlearnable} from partitioned data. As a result, there is a fundamental limitation to the achievable accuracy of prediction.


\subsection{Contributions}

A central contribution of this chapter is the design and analysis of the Partitioned Side Information Learning (PSIL) algorithm for combining the primary and secondary classifiers. The proposed algorithm first trains the secondary classifier and then uses it to impute the missing side information in the primary classifier’s training dataset. However, the estimated side information may not be fully accurate, as its accuracy depends on both the size of the secondary classifier’s training data and the complexity of the relationship between features and side information. Furthermore, since no historical data point contains both the true label and true side information (due to the partitioned dataset), any observed correlation between the label and the imputed side information could be spurious. To address this, we introduce a regularization mechanism that controls the influence of the estimated side information during training. The inverse of the regularization strength can be interpreted as a degree of trust in the imputed values. A high regularization value corresponds to little trust on the side information, while a low regularization value corresponds to greater trust. The optimal choice of this regularization parameter depends on knowledge of the true relationship between candidate aptitude, the label, and side information.

We establish a theoretical upper bound on the excess risk of the resulting composite classifier as a function of the regularization parameter. This upper bound exhibits the expected limiting behaviors: in the case of zero trust (infinite regularization), the bound reduces to that of the primary classifier alone; with complete trust (zero regularization), it becomes (roughly) the sum of the bounds for the primary and secondary classifiers. Importantly, our analysis shows that this upper bound on excess risk does not necessarily vanish, even with infinite training data. Indeed we confirm the non-vanishingness of excess risk by establishing a corresponding non-vanishing minimax lower bound on the excess risk. This fundamental limitation arises from the partitioned data setting, where features, labels, and side information are never jointly observed.


We do an ex-post analysis of the upper bound to find the optimal regularization assuming the data distribution is known. We find that the upper bound exhibits a conditional threshold-like behavior: depending on the underlying signal-imputation ratio, the bound-minimizing regularizer may take an extreme value—either fully trusting or completely ignoring the side information—or settle at an interior value. Though this conclusion pertains only to the upper bound and may not reflect the behavior of the actual excess risk, our simulations for the actual excess risk suggest that a similar behavior is observed there too.


\subsection{Related Work}

This chapter contributes to the literature on learning with auxiliary and incomplete information, bridging Learning Using Privileged Information (LUPI), missing data imputation, and transfer learning. This chapter is derived from our manuscript~\cite{singh2026side_journal}. To the best of our knowledge, the setting of training on partitioned datasets where side information is available only at test time has not been previously studied, and was first investigated in our work~\cite{singh2026side_journal}.

Learning Using Privileged Information (LUPI), introduced by \cite{vapnik2009new}, is a learning paradigm that leverages additional information available during training, but not at test time, to enhance learning performance. This is analogous to a teacher who, in addition to providing answers, also offers hints or explanations to students.
Subsequent works have extended the LUPI framework to various domains, including \cite{vapnik2015learning, sacco2021neural, song2021learning, wang2018learning}. In contrast, this chapter considers a setting where side information is not observed jointly with labels during training (i.e., the dataset is partitioned), and becomes available only at test time, making it distinct from the above approaches.

Our problem can be viewed as learning from an incomplete training dataset, in our case, a partitioned dataset. A rich body of literature addresses missing data imputation, including \cite{bertsimas2018predictive, van2007multiple, ghahramani1993supervised, stekhoven2012missforest, nguyen2023imputation}. Unlike these works, which primarily focus on the soundness of imputation methods, our approach treats imputation as an intermediate step, with the main objective of analyzing the excess risk of the proposed algorithm. Additionally, we apply regularization to the imputed values based on our confidence in their accuracy, introducing an added layer of caution. This emphasis on excess risk analysis distinguishes our approach from existing imputation-focused literature.

This chapter shares certain similarities with transfer learning, as discussed in \cite{hosna2022transfer}, where knowledge from an auxiliary task is used to improve performance on a main task. However, in standard transfer learning, auxiliary task labels are typically available during both training and testing, and are jointly observed with the main task labels during training. In contrast, our setting involves side information that is not jointly observed with labels during training, making it structurally different from conventional transfer learning approaches.

Several other works, such as \cite{adel2023general, chen2012boosting, jonschkowski2015patterns, mollaysa2017regularising}, define and utilize side information in their own ways across different contexts. A common aspect of these approaches is that side information is jointly observed with labels in the training data and is typically not available at test time (with the exception of \cite{adel2023general}, where side information is also available at test time). In contrast, our setting differs in formulation: we consider a partitioned training dataset and assume that side information is available only at test time.








